\section{Collision onsets and the information in their records}
\label{sec:geometric-results}
A maximal-collision event near its first possible time is supported on
an almost normal reflected segment. Its probability is exponentially
small in the number of flights, although the transverse sublevel has
only two free endpoint coordinates. The problem is to control that small
flux relatively, on a neighborhood which does not shrink with the
segment length. We prove such a control and pass from finite segments
to two semi-infinite boundary layers. The resulting nonlinear limiting
law distinguishes configurations whose entire leading statistical
hierarchies coincide.

There are three distinct levels of observation. The leading scalar
probability retains a determinant ratio. The leading conditional
endpoint metric retains the two contact curvatures separately, but all
its collision orders are already determined at one flight. The nonlinear
fixed-offset law retains higher boundary information. Theorem~\ref{thm:v4-jet-fiber}
separates the third level from both preceding ones by actual analytic
obstacles, with a nonzero coefficient computed in closed form. Thus
collision-order uniformity is used to establish a new limiting law,
not to count repeated leading coordinates as new inverse invariants.

\subsection{Periodic configurations and the observation}
Fix a full-rank lattice $\Lambda\subset\R^2$ and a positive finite number of bounded
$C^\infty$ strictly convex obstacles $D_1,\ldots,D_s$ with everywhere
positive curvature. All distinct translates $D_a+\lambda$, $\lambda\in
\Lambda$, have disjoint closures. Let $A>0$ be the free area per lattice
cell. The unit-speed billiard is prepared with normalized phase volume
$\dd q\dd\vartheta/(2\pi A)$ on the torus complement. We observe the
collision count $N_t$ and, when specified, the first and last collision
positions. Finite horizon, congruence of the obstacles, and symmetry
are not assumed.

Consider a fixed such configuration and a smooth finite-dimensional
family $D_{a,\xi}$ through it. The parameter $\xi$ ranges over a compact
set $K_0$ in a sufficiently small neighborhood of the reference parameter.
All conclusions are local about an arbitrary reference configuration,
not only about circles. A compact family is covered by finitely many
such neighborhoods, with common constants after taking finite extrema.
Boundaries may be represented by smooth support functions. At every
fixed derivative order, constants can depend on the corresponding
higher norms of this family, on its positive separation and curvature
bounds, and on the lattice. They do not depend on the collision order.

Let $\mathcal E$ be the finite set of oriented closest pairs attaining
the minimum gap in the reference configuration, modulo common lattice
translation. An index $e=(a,b,\lambda)$ represents a segment from
$D_{a,\xi}$ to $D_{b,\xi}+\lambda$; reversing it gives
$\bar e=(b,a,-\lambda)$. Section~\ref{sec:g-localization} proves that
these pairs continue smoothly, remain unobstructed, and include every
channel contributing to a common ground-onset collar. Write
\begin{equation}\label{eq:g-general-data}
 \begin{split}
 g_*&=\min_{e\in\mathcal E}g_e,\qquad
 c_{e,b}=1+g_e\kappa_{e,b}\quad(b=0,1),\\
 c_e&=\sqrt{c_{e,0}c_{e,1}},\qquad \gamma_e=\arcosh c_e,
 \qquad d_{j,e}=t-jg_e.
 \end{split}
\end{equation}
Here $\kappa_{e,0}$ and $\kappa_{e,1}$ are the positive curvatures at
its two contact points. They need not be equal. A channel event
$E_{j,e}(d)$ consists of the $j+1$ impacts alternating between the
specified facing patches in the window $t=jg_e+d$. These patches are
fixed around the reference contacts and chosen before observation.
They identify a measurable part of a physical collision record, not an
extra latent observation. At a nonminimal channel's own onset this
selection is required; the unselected count law there is a different
observation.

Parameter derivatives are multi-index derivatives in a fixed Euclidean
parameter chart; offset derivatives hold $d$ fixed as an independent
coordinate. Matrix norms are Euclidean operator norms unless a trace
norm is explicitly indicated. The growing record arrays later use the
maximum norm. These conventions distinguish uniform normalized
amplitudes from physical-time derivatives, which can carry powers of $j$.

\subsection{Uniform probability and conditional metric}
\begin{theorem}[Uniform thresholds for periodic dispersing configurations]
\label{thm:g-stability}
There is $\varepsilon_0>0$, independent of every integer $j\ge1$, such
that $\varepsilon_0<\min_{K_0}g_*/4$ and, for
$0\le t-jg_*<\varepsilon_0$,
\begin{equation}\label{eq:g-competition}
 \Pr_\xi(N_t\ge j+1)=\Pr_\xi(N_t=j+1)
 =\sum_{e\in\mathcal E}
 \frac{d_{j,e,+}^{\,2}}{2A(\xi)\sinh(j\gamma_e(\xi))}
       [1+d_{j,e,+}\mathcal R_{j,e}(\xi,d_{j,e,+})].
\end{equation}
Every fixed mixed derivative of $\mathcal R_{j,e}$ is bounded uniformly
in $j$; it has a smooth right extension at offset zero. For $t\le jg_*$
the probability is zero. At the selected channel's own window
$t=jg_e+d$, $0<d<\varepsilon_0$, the same individual summand is exactly
$\Pr_\xi(E_{j,e}(d))$, whether or not $e$ minimizes the gap.
\end{theorem}

The exponent two is independent of $j$, but its coefficient is
exponentially small. The essential estimate is relative to that
coefficient. The proof uses a weighted Dirichlet construction and a
trace-norm estimate for its relative determinant, rather than a
long initial-value iteration. The local hypotheses for this construction
are verified at every reference configuration above. In the triangular
one-obstacle near-circle class, $\mathcal E$ consists of the six unit
lattice directions, so \eqref{eq:g-competition} is its specialization, with all competing onsets.

For the next theorem, lift the first and last collision positions to
the fixed contact patches in $\R^2$, and denote them by $Q_0,Q_j$.
Define the two scalar conditional variances
\begin{equation}\label{eq:v3-observed-variance}
 V_{j,b}(d)=\frac1d\left\{
 \E(|Q_b|^2\mid E_{j,e}(d))-
       |\E(Q_b\mid E_{j,e}(d))|^2\right\},\qquad b=0,j.
\end{equation}
These are translation- and rotation-invariant within the chosen lifts.
Their definition requires neither the unknown contact point nor its
tangent or curvature. The gap $g_e$ is read from the channel onset.

\begin{theorem}[Uniform two-variance curvature recovery]
\label{thm:v3-tomography}
For every odd $j\ge1$ the right limits $v_{j,0},v_{j,1}$ of
$V_{j,0}(d),V_{j,j}(d)$ exist, and
\begin{equation}\label{eq:v3-variance-limits}
 (v_{j,0},v_{j,1})=
 \frac{g_e\coth(j\gamma_e)}{3\sinh\gamma_e}
 \left(\sqrt{\frac{c_{e,1}}{c_{e,0}}},
       \sqrt{\frac{c_{e,0}}{c_{e,1}}}\right).
\end{equation}
Each difference $V_{j,b}(d)-v_{j,b}$ is $d$ times a smooth function
with uniform fixed-order derivative bounds. At known $g_e$, these two
limits determine the ordered pair $(\kappa_{e,0},\kappa_{e,1})$ without
knowing $A$. The inverse is given by
\begin{equation}\label{eq:v3-inverse-map}
 \begin{gathered}
 \frac{\coth(j\gamma)}{\sinh\gamma}
        =\frac3g\sqrt{v_0v_1},\qquad
 r=\sqrt{v_1/v_0},\\
 \kappa_0=\frac{r\cosh\gamma-1}{g},\qquad
 \kappa_1=\frac{r^{-1}\cosh\gamma-1}{g}.
 \end{gathered}
\end{equation}
On every compact positive gap--curvature box this observation map and
its inverse are Lipschitz, with constants independent of odd $j$.
They remain uniformly conditioned at $\kappa_0=\kappa_1$.
If both measured variances have error at most $\delta$, a reconstruction
constrained to such a box has curvature error at most $C(d+\delta)$.
\end{theorem}

Thus the full conditional metric retains information lost by the
scalar flux coefficient. Theorem~\ref{thm:v3-fibers} realizes the
scalar ambiguity by actual smooth periodic obstacles with the same gap,
area and every leading count amplitude, while their endpoint curvatures
vary. The two variances distinguish that family. The statement is local
contact-curvature recovery, not global shape rigidity. Conditioning on
an exponentially rare event is not free: Proposition~\ref{prop:v3-sampling}
records the number of independent physical preparations required.

\subsection{Sources and retained consequences}
Let $x_{e,0},x_{e,1}$ be the normal outgoing states and $f_\xi$ a smooth
collision mark near them. Put $a_{e,b}=f_\xi(x_{e,b})$,
$\bar f_e=(a_{e,0}+a_{e,1})/2$, and
$w_{j,e}=\sum_{i=0}^j a_{e,i\bmod2}$.
\begin{theorem}[Marked response and the exact axial source domain]
\label{thm:g-marked}
At the selected channel's own threshold,
\begin{equation}\label{eq:g-marked}
 Z_{j,e}(\xi,q,d)=
 \frac{d^2e^{qw_{j,e}(\xi)}}{2A(\xi)\sinh(j\gamma_e(\xi))}
       [1+d\mathcal R^f_{j,e}(\xi,q,d)].
\end{equation}
Every fixed mixed derivative of the remainder is bounded uniformly in
$j$ on compact parameter and source sets. The complete marked event in
the ground collar is the sum at its positive channel offsets. For fixed
$\xi$, real $q$, and sufficiently small fixed $d\ge0$, the series
$\sum_{j,e}Z_{j,e}(\xi,q,d)/d^2$ converges exactly when
\begin{equation}\label{eq:g-source-domain}
           \max_{e\in\mathcal E}(q\bar f_e-\gamma_e)<0.
\end{equation}
Right limits are used at $d=0$. Every fixed mixed derivative is summable
on compact subsets of this domain. Finitely many real sources are
allowed, with the scalar product in place of $q\bar f_e$.
\end{theorem}

This series ranges over different onset windows, not over all
itineraries in one unrestricted long window. The source formula,
interface jumps and record response are consequences of the same
uniform construction, not separate claims of a general limit theory.
The exact equal-gap inverse theorem is retained in
Section~\ref{sec:g-inverse}. It uses central symmetry to remove the
two-endpoint ambiguity and equal gaps to retain all channels at every
collision order. The recovered area is a symmetric function of the
three recovered curvatures in that chosen family, not a fourth
independent unknown. Section~\ref{sec:v3-amplitude-data} proves this
identity and the loss of first-order sensitivity at coalescence.

\subsection{Nonlinear long bridges}
For one selected channel set
\[
 E_j(u,v)=W_{j,e}(u,v)-jg_e,\qquad
 b_j(u,v)=\frac{-W_{j,e,uv}(u,v)}{-W_{j,e,uv}(0,0)},\qquad
 F_j(d)=\frac{2A\sinh(j\gamma_e)}{d^2}\Pr(E_{j,e}(d)).
\]
Right limits are used at zero. We construct positive smooth amplitudes
$B_b$ and strictly convex boundary actions $S_b$, $b=0,1$, from
semi-infinite reflected segments. Theorem~\ref{thm:v4-factorization}
proves, for $p=j\bmod2$, exponential convergence
\[
 E_j\longrightarrow S_0(u)+S_p(v),\qquad
 b_j\longrightarrow B_0(u)B_p(v)
\]
in every fixed differentiability norm on a common endpoint neighborhood.
The second convergence is relative to the exponentially small twist.
Theorem~\ref{thm:v4-law} gives the corresponding smooth convergence of
$F_j(d)$ on a common closed offset interval. Corollary~\ref{cor:v4-poles}
then identifies the first singularities of the generating function of
these actual, finite-offset probabilities.

The analytic family in Theorem~\ref{thm:v4-jet-fiber} has fixed gap,
area, contact curvatures, and every leading endpoint matrix. Its limiting
coefficient $\partial_d\mathcal F_p(0)$ nevertheless changes, with derivative
$\sqrt3/2$ at the circular member in the normalization specified there.
This is a concrete consequence of the nonlinear boundary construction.
It neither asserts global shape rigidity nor uses a characteristic
function for unrestricted itineraries.

Sections~\ref{sec:g-localization}--\ref{sec:g-integration} establish
the geometric localization, common Dirichlet chart, and physical
integration. Sections~\ref{sec:v4-boundary}--\ref{sec:v4-nonlinear}
prove the nonlinear limiting statements and their geometric separation.
The conditional inverse is followed by its leading-data factorization
and a fixed-order extrapolation theorem with acquisition and timing
bounds. The exact unlabelled inverse is supplemented by a sharp
circular-reference stability estimate, not a pairwise noisy inverse for
arbitrary triples. Smooth and moving-cut record response, the circular
appendices, and the unchanged fixed-window companion complete the
article.
